\slajdid{02-s054}
\begin{frame}[label=02-s054]
\gusttekst
Potpuno identična situacija nastaje ako je funkcija prikazana u obliku proizvoda zbirova, odnosno minimizacija urađena prekrivanjem logičkih nula. Razlika je u tome što će se prilikom prelaska sa površine na površinu koje obezbeđuju logičke nule, pojaviti lažna jedinica. Rešenje je isto dodavanje zajedničke površine koja će povezati susedne površine koje nemaju zajedničkih polja.

Liči da smo na lak način razrešili problem, međutim u praksi je problem složeniji, pošto je moguće da se više ulaznih signala istovremeno menja. Osnovni problem koji nastaje jeste kada na kombinacionu mrežu gledajući samo I i ILI deo signali dolaze sa različitim kašnjenjem. Jedino moguće rešenje jeste da se ta kašnjenja ujednačavaju, dodavanjem dodatnih komponenti, invertora, ili da se preuređuje mreža i odstupa (kao što već jesmo) od „minimalne realizacije“ a što se u većini slučajeva radi heuristički ili simulacija. Suštinski presudno je razumevanje i iskustvo od strane projektanta.
\end{frame}
